\section{Stage III: The Safepoint Tax}
\label{sec:polls}

My guide suggested pricing the poll with two experiments that bracket its cost.

\subsection{Two Ways to Remove a Poll}
\begin{itemize}
  \item \textbf{NOP patch (lower bound).} \code{tools/nop_polls.py}
        overwrites every \ainl|cmp (%r14),%r15| plus \code{jbe} pair (a poll)
        in the linked binary with multi-byte NOPs, aborting unless every poll
        is gone and the count of similar-looking allocation checks is
        unchanged (in the first binary patched, 853 polls and 1,263 checks).
        The output stays bit-identical, and registers and layout are
        untouched.
  \item \textbf{Compiler flag (upper bound).} \code{-disable-poll-insertion}
        recompiles without polls, so the compiler can re-optimise around
        their absence.
\end{itemize}

Removing polls is correct for single-domain microgpt, not in general
(Section~\ref{sec:reality}).

\subsection{Results: What You Remove Is Not What You Pay}
\label{fig:polls}
On the release's plain kernel, which executes about four times as many
back-edge polls as the unrolled one, the NOP patch gains
@P{oxcaml_unsafe_plain,oxcaml_unsafe_plain_pollnop}@ and the flag
@P{oxcaml_unsafe_plain,oxcaml_unsafe_plain_pollfree}@. On the unrolled kernel
the NOP patch gains @P{oxcaml_unsafe_unroll,oxcaml_unsafe_unroll_pollnop}@
(interquartile range @I{oxcaml_unsafe_unroll,oxcaml_unsafe_unroll_pollnop}@)
and the flag @P{oxcaml_unsafe_unroll,oxcaml_unsafe_unroll_pollfree}@.
Executing the poll costs a few percent at most; recompiling without polls can
be worth several times that. On the unrolled kernel the two are close because
5.4.0-ox7's poll-free build re-tags the index with an extra \code{lea} that
forces a spill: it is @A{oxcaml52_unsafe_unroll_pollfree,oxcaml_unsafe_unroll_pollfree}@
slower than the same build on the previous release, 5.2.0+ox.

\subsection{The \texorpdfstring{\%r11}{r11} Exclusion}
\label{sec:r11}
On 5.2.0+ox, without polls \code{linear} shrinks from 102 to 91 instructions, its one
stack spill disappears, and register \reg{r11} goes from 0 uses to 8. The
backend declares \reg{r11} destroyed at every poll and allocation
(\code{destroyed_at_alloc_or_poll} in \code{backend/amd64/proc.ml}), so with a
poll in the loop the allocator spills a row base.

A binary patch that rewrote the two spill instructions into moves through
\reg{r11} (same instruction count, no spill) gained only $+0.3\%$ (7/11
rounds, a separate run): the slot is in L1 and off the critical floating-point
chain, so the spill is not where the poll's cost lies. The runtime does save
\reg{r11} in \code{SAVE_ALL_REGS}~\cite{ocaml_runtime}, and an OCaml core
developer told me the exclusion may be redundant in the non-PLT
configuration; settling it needs a test that forces the slow path with a
pointer in \reg{r11}.

What, then, makes the flag worth more than the poll it removes? On the
release it changes code layout, block structure and register allocation
together, and these experiments cannot separate them. On the patched compiler
most of the value of removing polls is in the optimisations their absence
\emph{permits} (Section~\ref{sec:reality}).

\paragraph{Configuration for the next stages.} Stages IV to VI develop the
passes on compilers built without poll insertion, to study the address
arithmetic with no poll in the loop; Stage VII measures them with polls on,
the configuration OxCaml is released with. The final measurement of each pass
(Table~\ref{tab:passes}) uses the pull-request compiler configured like the
release, with counted loops (Section~\ref{sec:counted}).
